% !TeX root = ../../Book3.tex
% 纲要标题：### 9.4 局部维数、参数系与嵌入维数
% 纲要位置：工作记录/计划纲要.md，第 2097 行。
% 纲要细目（写作范围）：
% * 参数系
% * 嵌入维数与 Krull 维数
% * 正则局部环的定义先行，结构理论留给第18章
% ---


\section{局部维数、参数系与嵌入维数}
\label{b3:sec:0904}

在 Noether 局部环中，维数有两种容易混淆的比较对象：使零点只剩闭点需要多少个方程，以及生成整个极大理想需要多少个元素。前者等于 Krull 维数，后者可能更大。两者相等时得到正则局部环的定义。


\subsection{参数系}
\label{b3:subsec:0904:01}

\begin{definition}\label{b3:def:system-of-parameters}
设 \((R,\mathfrak m)\ne0\) 为交换 Noether 局部环，\(d=\dim R\)。若一组元素
\(x_1,\ldots,x_d\in\mathfrak m\) 满足
\(\sqrt{(x_1,\ldots,x_d)}=\mathfrak m\)，则称它们为 \(R\) 的\term[canshuxi]{参数系}[system of parameters]，称它们生成的理想为\term[canshu-lixiang]{参数理想}[parameter ideal]。
\end{definition}

\begin{theorem}\label{b3:thm:system-of-parameters}
设 \((R,\mathfrak m)\ne0\) 为交换 Noether 局部环，\(d=\dim R\)。则 \(R\) 有参数系，且 \(d\) 等于生成一个根为 \(\mathfrak m\) 的理想所需的最少元素个数。
若 \(x_1,\ldots,x_d\) 是参数系，则对 \(0\le i\le d\)，有
\[
\dim R/(x_1,\ldots,x_i)=d-i.
\]
此外，对任意 \(a_1,\ldots,a_r\in\mathfrak m\)，都有
\(\dim R/(a_1,\ldots,a_r)\ge d-r\)。
\end{theorem}
\begin{proof}
存在性已由\cref{b3:lem:dimension-cutting}证明，高度定理说明少于 \(d\) 个元素不可能生成根为 \(\mathfrak m\) 的理想。
设 \(e=\dim R/(a_1,\ldots,a_r)\)，在此商环中选一个有 \(e\) 个元素的参数系并提升。
它与 \(a_1,\ldots,a_r\) 合起来生成根为 \(\mathfrak m\) 的理想，所以 \(d\le r+e\)。
对参数系的前 \(i\) 项应用这个不等式，得到商维数至少是 \(d-i\)；剩余 \(d-i\) 项在商环中生成极大根理想，高度定理又给出反向不等式。
\end{proof}

\(R/(x_1,\ldots,x_d)\) 是零维 Noether 局部环，故是 Artin 环，且作为 \(R\) 模具有有限长度。
反之，若商模具有有限长度，其支集只含闭点，于是生成理想的根为 \(\mathfrak m\)。所以参数系也可以通过商环的有限长度识别。

\ProseParagraphBreak
例如在 \(k[x,y]_{(x,y)}\) 中，\(x,y\) 和 \(x^2,y^3\) 都是参数系，后者并不生成极大理想。
在尖点局部环 \(k[t^2,t^3]_{(t^2,t^3)}\) 中，\(t^2\) 单独构成参数系：模掉它后
\((t^3)^2=(t^2)^3=0\)，故商环只有一个素理想。参数并不要求能作自由的坐标。

\subsection{嵌入维数与 Krull 维数}
\label{b3:subsec:0904:02}

\begin{definition}\label{b3:def:embedding-dimension}
设 \((R,\mathfrak m)\ne0\) 为交换 Noether 局部环，剩余域为 \(\kappa=R/\mathfrak m\)。向量空间 \(\mathfrak m/\mathfrak m^2\) 在 \(\kappa\) 上的维数
\[
\dim_\kappa(\mathfrak m/\mathfrak m^2)
\]
称为 \(R\) 的\term[qianru-weishu]{嵌入维数}[embedding dimension]，记为 \(\operatorname{edim}R\)。
\end{definition}
\symidx{\(\operatorname{edim}R\)：局部环的嵌入维数}

\(\mathfrak m/\mathfrak m^2\) 上的标量作用来自 \(R\)，并因 \(\mathfrak m\) 零化该商模而降为 \(\kappa\) 作用。
\cref{b3:cor:local-generators}说明，嵌入维数恰好等于极大理想的最少生成元数。
由高度定理立即得到
\begin{equation}\label{b3:eq:dimension-embedding-bound}
\dim R\le\operatorname{edim}R.
\end{equation}

这个向量空间仅保留函数的一次信息；乘积落入 \(\mathfrak m^2\) 而消失。取
\[
R=k[x_1,\ldots,x_n]_{(x_1,\ldots,x_n)},
\]
则变量的像构成
\(\mathfrak m/\mathfrak m^2\) 的基。生成性显然；若线性式属于局部化后的平方理想，则乘一个常数项非零的分母后，其一次部分仍是该线性式乘以非零常数，故各系数必须为零。因此嵌入维数为 \(n\)。坐标素链在此局部化中全部保留，Krull 维数也为 \(n\)。

\ProseParagraphBreak
尖点的情形不同。\(t^2,t^3\) 在 \(\mathfrak m/\mathfrak m^2\) 中线性独立，因为平方理想中每个非零元素的最低次数至少为四，局部化的单位分母不改变这一结论。因此该环的嵌入维数为二，Krull 维数为一。这两个数字之差已经能辨认这个点与直线上的点不同。

\subsection{正则局部环的定义}
\label{b3:subsec:0904:03}

\begin{definition}\label{b3:def:regular-local-ring}
若非零交换 Noether 局部环 \((R,\mathfrak m)\) 满足
\(\dim R=\dim_{R/\mathfrak m}(\mathfrak m/\mathfrak m^2)\)，则称它为
\term[zhengze-jubuhuan]{正则局部环}[regular local ring]。
\end{definition}

域是零维正则局部环；反过来，零维正则局部环的极大理想满足
\(\mathfrak m/\mathfrak m^2=0\)，由 Nakayama 引理有 \(\mathfrak m=0\)，所以它是域。
上小节的多项式局部环是正则的，尖点局部环则不是。
另一个简单的非正则例子为 \(k[\epsilon]/(\epsilon^2)\)：其 Krull 维数为零，而嵌入维数为一。

\ProseParagraphBreak
设 \((R,\mathfrak m)\) 为 \(d\) 维正则局部环。极大理想 \(\mathfrak m\) 的任一极小生成组
都有 \(d\) 项，并且生成的理想就是 \(\mathfrak m\)，因而构成参数系；这样的生成组称为\term[zhengze-canshuxi]{正则参数系}[regular system of parameters]。
正则序列在\chapref{b3:ch:15}引入，正则参数系构成正则序列的证明见\secref{b3:sec:1801}。
正规表示整闭，正则表示上述维数等式，它们是不同定义。下一章将在一维整环情形研究二者与离散赋值环的关系。

\begin{table}[htbp]
\centering
\caption[维数类不变量]{维数类不变量。\(R\) 为交换含幺环，\(M\) 为 \(R\) 模；局部环行中的 \(\mathfrak m\) 表示唯一极大理想，\(\kappa=R/\mathfrak m\)。空支集与零环的维数取 \(-\infty\)。}
\label{b3:tab:dimension-invariants}
\begin{tabularx}{\textwidth}{@{}>{\raggedright\arraybackslash}p{0.19\textwidth}>{\raggedright\arraybackslash}p{0.23\textwidth}>{\raggedright\arraybackslash}X@{}}
\toprule
不变量 & 条件 & 所测量的对象与结论\\
\midrule
\(\operatorname{ht}\mathfrak p\) & \(\mathfrak p\in\operatorname{Spec}R\) & 以 \(\mathfrak p\) 为末项的严格素链长的上确界，等于 \(\dim R_{\mathfrak p}\)\\
\(\dim(R/\mathfrak p)\) & \(\mathfrak p\in\operatorname{Spec}R\) & 以 \(\mathfrak p\) 为首项的严格素链长的上确界，即闭子空间 \(V(\mathfrak p)\) 的维数\\
\(\dim_R M\) & 任意 \(R\) 模 \(M\) & 支集 \(\operatorname{Supp}_R M\) 内的严格素链长的上确界；若 \(M\) 有限生成，等于 \(\dim(R/\operatorname{Ann}_R M)\)\\
\(\operatorname{edim}R\) & \((R,\mathfrak m)\ne0\) 为 Noether 局部环 & \(\dim_\kappa(\mathfrak m/\mathfrak m^2)\)，等于生成 \(\mathfrak m\) 所需的最少元素个数\\
生成 \(\mathfrak m\) 准素理想的最少元素数 & \((R,\mathfrak m)\ne0\) 为 Noether 局部环 & 在所有 \(\mathfrak m\) 准素理想的生成组中取元素个数的最小值，等于 \(\dim R\)；达到最小值的生成组是参数系\\
\bottomrule
\end{tabularx}
\end{table}

\cref{b3:tab:dimension-invariants}中的高度测量指定素理想以下的链，商环维数测量它以上的链。若各量有限，\cref{b3:prop:dimension-basic-comparison}给出
\[
\operatorname{ht}\mathfrak p+\dim(R/\mathfrak p)\le\dim R.
\]
\begin{samepage}
取一条经过 \(\mathfrak p\)、以极大理想 \(\mathfrak m\) 为末项的素链，可以在 \(\mathfrak p\) 处分成两段：
\[
\underbrace{\mathfrak p_0\subsetneq\cdots\subsetneq\mathfrak p}
_{u\le\operatorname{ht}\mathfrak p}
\underbrace{\subsetneq\cdots\subsetneq\mathfrak m}
_{v\le\dim(R/\mathfrak p)}.
\]
这里 \(u,v\) 是两段各自的长度。对两段分别取最长链再拼接，就得到上述不等式；原环的最长链未必经过这个指定素理想。
\end{samepage}
